\begin{officialform}[BulletinVI]
\FormHeading
\FormLine{Mã số đề tài: \ProjectCode}
\FormLine{Tên đề tài: \ProjectTitle}
\FormLine{Thời gian thực hiện: \ProjectPeriod}
\FormLine{Tổng kinh phí: \ProjectBudget}
\FormLine{Chủ nhiệm đề tài: \PrincipalName\ (\PrincipalPhone)}
\FormLine{Thành viên tham gia nghiên cứu: \ProjectMembers}
\FormLine{Cán bộ hướng dẫn: \SupervisorName\ (\SupervisorEmail)}
% Nêu lý do cần nghiên cứu.
\FormLine{Tính cấp thiết: Nhu cầu dùng chung ứng dụng Kanban cho nhiều lớp học và nhóm nghiên cứu đặt ra yêu cầu định danh, cô lập dữ liệu và duy trì nhất quán trạng thái công việc.}
% Nêu mục tiêu đề tài.
\FormLine{Mục tiêu: Thiết kế, xây dựng và đánh giá kiến trúc đa thuê bao cho ứng dụng quản lý công việc trong môi trường đại học.}
% Nêu phương pháp chính.
\FormLine{Phương pháp nghiên cứu: Tổng hợp lý thuyết, phân tích đánh đổi kiến trúc, xây dựng ứng dụng và thực nghiệm kiểm chứng.}
% Nêu các nội dung đã thực hiện.
\FormLine{Nội dung nghiên cứu: Định danh thuê bao; cô lập dữ liệu và phân quyền; nhất quán công việc; thiết kế, triển khai và đánh giá ứng dụng.}
% Tóm tắt kết quả có căn cứ.
\FormLine{Kết quả đạt được: \PendingField{Kết quả đạt được}}
% Nêu đóng góp của đề tài.
\FormLine{Ý nghĩa: \PendingField{Ý nghĩa}}
% Nêu điều kiện và phạm vi áp dụng.
\FormLine{Khả năng ứng dụng: \PendingField{Khả năng ứng dụng}}
\end{officialform}
